%------------------------------------------------------------------------------%
\begin{question}
  Calcule quanto
  \[
    f(x, y, z) = y\sin(x) + 2yz
  \]
  irá variar se nos movermos 0,1 unidades a partir de \((0, 1, 0)\) \\
  na direção do ponto \((2, 2, -2)\)
\end{question}

%------------------------------------------------------------------------------%
\begin{resolution}{Direção do movimento}
  \[
    \pause
    v
    \pause
    \;=\; \vetortri{2}{2}{-2}
        - \vetortri{0}{1}{ 0}
    \pause
    \;=\; \vetortri{2}{1}{-2}
    \pause
    \;=\; 2\veci + \vecj -2\veck
  \]
  \pause
  Normalizando o vetor
  \[
    \pause
    u
    \pause
    = \frac{v}{\abs{v}}
    \pause
    = \frac{2\veci + \vecj -2\veck}{\sqrt{2^2 + 1^2 + (-2)^2}}
    \pause
    = \frac{2\veci + \vecj -2\veck}{\sqrt{9}}
    \pause
    = \frac{2}{3} \veci + \frac{1}{3} \vecj - \frac{2}{3} \veck
  \]
\end{resolution}

%------------------------------------------------------------------------------%
\begin{resolution}{Derivadas parciais}
  \begin{align*}
    \onslide<+->{\D{f}{x}}
    \onslide<+->{& = \D{}{x}\left(y\sin(x) + 2yz\right)}
    \onslide<+->{ = y\cos(x) \\[5mm]}
    \onslide<+->{\D{f}{y} & = \D{}{y}\left(y\sin(x) + 2yz\right)}
    \onslide<+->{ = \sin(x) + 2z\\[5mm]}
    \onslide<+->{\D{f}{z} & = \D{}{z}\left(y\sin(x) + 2yz\right)}
    \onslide<+->{ = 2y}
  \end{align*}
\end{resolution}

%------------------------------------------------------------------------------%
\begin{resolution}{Gradiente de $f$}

  \[
    \onslide<+->{\nabla f}
    \onslide<+->{= (y\cos(x))\veci + (\sin(x) + 2z)\vecj + (2y)\veck}
    \onslide<+->{= \vetortri{y\cos(x)}{\sin(x) + 2z}{2y}}
  \]

  \[
    \onslide<+->{\nabla f(0, 1, 0)}
    \onslide<+->{= (1\cos(0))\veci + (\sin(0) + 2\times 0)\vecj + (2\times 1)\veck}
    \onslide<+->{= \veci + 2\veck}
    \onslide<+->{= \vetortri{1}{0}{2}}
  \]
\end{resolution}

%------------------------------------------------------------------------------%
\begin{resolution}{Derivada direcional}
  \[
    D_u f(0, 1, 0)
    \pause
    \;=\; \nabla f(0, 1, 0) \cdot u
  \]

  \[
    \hphantom{D_u f(0, 1, 0)}
    \pause
    \;=\; \left(\veci + 2\veck\right)
    \cdot
    \left(\frac{2}{3} \veci + \frac{1}{3} \vecj - \frac{2}{3} \veck\right)
  \]

  \[
    \hphantom{D_u f(0, 1, 0)}
    \pause
    \;=\; 1\times\frac{2}{3} + 0\times\frac{1}{3} + 2 \times\left(- \frac{2}{3}\right)
  \]

  \[
    \hphantom{D_u f(0, 1, 0)}
    \pause
    \;=\; \frac{2}{3} - \frac{4}{3}
    \pause
    \;=\; -\frac{2}{3}
  \]
\end{resolution}

%------------------------------------------------------------------------------%
\begin{resolution}{Variação}
  \[
    \pause
    df
    \pause
    \;=\; D_u f(0, 1, 0) ds
  \]

  \[
    \pause
    \hphantom{df}
    \;=\; -\frac{2}{3}\times 0,1
    \pause
    \;=\; -\frac{2}{3} \frac{1}{10}
    \pause
    \;=\; -\frac{2}{30}
    \pause
    \;=\; -\frac{1}{15}
  \]

  \[
    \pause
    \hphantom{df}
    \approx -0,067
  \]
\end{resolution}

%------------------------------------------------------------------------------%